\subsection{The actual arithmetic correction space and its integrated automorphisms}

\begin{statement}{D.53}{Proved}{Exact local and transgression maps for the arithmetic moving corrections}
Return to the actual arithmetic groups of D.49--D.51:
\[
 \begin{gathered}
 G=G_3,\qquad n=576\cdot3^{97},\qquad
 K_3=\ker(\widehat B\to G),\\
 U=K_3^{\rm ab,(2)},\qquad V=U/2U,\qquad
 \mathscr W=(1-e_A)U,\qquad W=\mathscr W/2\mathscr W,\\
 q=|A|\ge3,\qquad
 w=\dim_{\mathbb F_2}W={n\over6}(1-q^{-1}).
 \end{gathered}
 \tag{D.53.1}
\]
Here \(A\) is the same actual normal elementary
three-subgroup generated by \(f_h\bmod G_3\).
For \(g=s,r,t\), write \(W^g=\ker(1-g:W\to W)\).

The complete allowable modulo-two conjugator-pair
space is
\[
 \mathcal P=
 \{(b,c)\in W\oplus W:(1-s)b+s(1-r)c=0\},
 \qquad \dim\mathcal P=w .
 \tag{D.53.2}
\]
Every pair induces the actual continuous derivation
\[
 \delta_{b,c}(s)=(1-s)b,\qquad
 \delta_{b,c}(r)=(1-r)c,\qquad
 \delta_{b,c}(t)=0 .
 \tag{D.53.3}
\]
Restriction to \(K_3\) gives a specified linear map
\[
       \mathcal P\longrightarrow\operatorname{Hom}_G(V,W).
 \tag{D.53.4}
\]
Its kernel is exactly the injective image
\[
 \begin{aligned}
 W^s\oplus W^r\oplus W^t&\longrightarrow\mathcal P,\\
 (u_s,u_r,d)&\longmapsto(u_s+d,u_r+d).
 \end{aligned}
 \tag{D.53.5}
\]
Thus the actual image \(\mathscr D\) has dimension
\[
       d=w-\dim W^s-\dim W^r-\dim W^t .
 \tag{D.53.6}
\]
The local and transgression restrictions are
computed by the exact sequences
\[
 \begin{gathered}
 H^i(G,W)=0\quad(i\ge0),\qquad
 H^1(\widehat B,W)\xrightarrow{\sim}\operatorname{Hom}_G(V,W),\\
 0\longrightarrow W/(W^s+W^r)
  \longrightarrow\operatorname{Hom}_G(V,W)
  \longrightarrow H^1(\langle s\rangle,W)
  \longrightarrow0,\\
 0\longrightarrow\mathscr D
  \longrightarrow W/(W^s+W^r)
  \xrightarrow{\ \kappa\ }W/(1-t)W
  \longrightarrow0,\\
 \kappa([v])=[s(1-r)v],\qquad
 H^1(\langle s\rangle,W)=\ker(1+s)/(1-s)W .
 \end{gathered}
 \tag{D.53.7}
\]
In particular vanishing transgression does not
remove the displayed local and cusp restrictions.

These are nonzero correction directions for the
actual arithmetic source. Their dimension satisfies
the explicit positive bound
\[
       d\ \ge\
          80\cdot3^{95}-4\cdot3^{45}-4\cdot3^{34}
          \ >0 .
 \tag{D.53.8}
\]
The proof uses the actual \(576\)-state marked
graph and its rank-\(97\) three-homology extension.
Its exact centralizer and Jordan data are
\[
 \begin{gathered}
 |C_Q(s)|=16,\qquad |C_Q(r)|=18,\\
 \dim(N^{\rm ab}/3)^s=45,\qquad
 (N^{\rm ab}/3)|_{\langle r\rangle}
      =J_3^{\,29}\oplus J_2^{\,5},\\
 |C_G(s)|=16\cdot3^{45},\qquad |C_G(r)|=3^{35}.
 \end{gathered}
 \tag{D.53.9}
\]
Here \(J_j\) is the length-\(j\) unipotent block
over \(\mathbb F_3\).

The correction directions integrate on the
complete actual D.50 source. Define
\[
 \mathscr P=
 \ker\bigl(\mathscr W\oplus\mathscr W
       \xrightarrow{\, (1-s)b+s(1-r)c\,}\mathscr W\bigr).
 \tag{D.53.10}
\]
This is a free \(\mathbb Z_2\)-module of rank \(w\);
its reduction modulo two maps onto \(\mathcal P\).
For the chosen D.50 automorphism \(\alpha_0\),
every pair \(2(b,c)\), with \((b,c)\in\mathscr P\),
gives an actual continuous peripheral automorphism
\(\alpha_{2(b,c)}\) of \(\mathcal Z_h\).
It fixes the cusp, retains the entire D.49 action
and parameter, and has the same action modulo two
as \(\alpha_0\).

The closed subgroup generated by
\(\alpha_{2(b,c)}\alpha_0^{-1}\) has a continuous
elementary abelian quotient of dimension \(d\),
detected by its homology action modulo four:
\[
 \alpha\longmapsto
       {\,\alpha|_{U/4U}-1\,\over2}\bmod2 .
 \tag{D.53.11}
\]
This is an integrated automorphism quotient of
the specified actual D.50 source. It is not a
claim that these automorphisms lift to
\(\widehat B\) or solve the remaining invariant
class-two equation.

The separate order-\(1080\) source provides an
evaluated check of the local ambiguities.
There
\[
 \begin{gathered}
 \dim W^s=64,\qquad \dim W^r=38,\qquad \dim W^t=0,\\
 \dim\mathcal P=120,\qquad \dim\ker(\mathcal P\to\mathscr D)=102,
 \qquad \dim\mathscr D=18 .
 \end{gathered}
 \tag{D.53.12}
\]
The same-action kernel has rank \(4\) in
D.52's \(13\)-coordinate invariant obstruction.
It repairs the previously unsuccessful
pairing-preserving action while keeping that
entire modulo-two homology matrix unchanged.
Its full corrected cusp word has
\(4\,220\,100\) letters and is checked directly
in the unquotiented class-two target; generation
has rank \(181\), and the homology deviation
still has rank \(16\). This finite check does
not evaluate the arithmetic invariant obstruction.
The manuscript passages are Part D(e),
section 51.16, folio 734 verso, scan 174,
and section 51.15, folios 725--727,
scans 156--160.
\end{statement}
\begin{proof}
\emph{The actual moving block and the transgression.}
The averaging element \(e_A\) is a central
idempotent of \(\mathbb F_2[G]\), since \(A\)
is normal and has odd order. It acts as
identity on the trivial module and as zero
on \(W\). A projective resolution of the
trivial module can be taken in the \(e_A\)
block: apply that idempotent to a free
resolution. Every homomorphism from that
resolution to \(W\) is zero. This proves
\(H^i(G,W)=0\) in every degree, including
degree zero.

The actual inflation--restriction sequence
for \(K_3\subset\widehat B\) now gives the
isomorphism in (D.53.7). Since the action
on \(W\) factors through \(G\), restriction
on \(K_3\) factors through its elementary
two-homology \(V\). Thus every
\(G\)-equivariant homomorphism has zero
transgression and represents a unique
cohomology class of continuous derivations.
The peripheral restrictions are computed next.

\emph{The edge map and the exact local restriction.}
Use the actual graph sequence
\[
 0\longrightarrow V\longrightarrow
 \mathbb F_2[G]\longrightarrow
 \mathbb F_2[G/\langle s\rangle]\oplus
 \mathbb F_2[G/\langle r\rangle]
 \longrightarrow\mathbb F_2\longrightarrow0 .
 \tag{D.53.13}
\]
Its boundary sends the identity edge to its
two incident vertices. Apply
\(\operatorname{Hom}_G(-,W)\), separating the
image of the boundary into a short exact
sequence. The just-proved vanishing of
all \(\operatorname{Ext}_G^i(\mathbb F_2,W)\)
removes the augmentation terms.
The edge module is projective, and the
\(r\)-vertex module is projective because
three is invertible in the coefficient
field. The remaining vertex extension
group is \(H^1(\langle s\rangle,W)\).
This gives the first short exact sequence
in (D.53.7).

More explicitly, its map from \(W\) is
the restriction to the cycle module of
the edge map \(e_g\mapsto gv\).
Its kernel is \(W^s+W^r\), the two vertex
values. Their intersection is \(W^G=0\).
The corresponding path derivation has
values
\[
                 \delta_v(s)=0,\qquad
                 \delta_v(r)=(1-r)v .
\]
It is locally a coboundary at both finite
peripheral subgroups. Conversely all such
cohomology classes lie in this edge-map
image, by the exact sequence. This proves
the claimed local restriction without
assuming that every abstract equivariant
map is peripheral.

The cusp value of this derivation is
\(s(1-r)v\). Changing it by an inner
derivation changes that value by
\((1-t)W\). Thus the remaining map is
exactly \(\kappa\). It is well-defined
on \(W/(W^s+W^r)\): an \(r\)-fixed
vector contributes zero, and an \(s\)-fixed
vector contributes \((1-t)v\).
It is onto because
\[
       (1-t)W+s(1-r)W
           =(1-s)W+s(1-r)W=W .
 \tag{D.53.14}
\]
The last equality is the augmentation
identity \(W_G=0\), already implied by
\(W_A=0\). Its kernel consists exactly
of the cusp-fixed peripheral restrictions.
This proves the second short exact sequence.

\emph{Conjugator parameters and their full kernel.}
The same augmentation identity proves
that the map defining \(\mathcal P\)
is onto. Its kernel therefore has
dimension \(w\). Each pair gives
(D.53.3); the two torsion relations
hold by telescoping, and the defining
equation makes its cusp value zero.

Its restriction to \(K_3\) is zero
exactly when the derivation inflates
from \(G\). Vanishing \(H^1(G,W)\)
makes it the inner derivation
\(g\mapsto(1-g)d\); \(d\) is unique
because \(W^G=0\).
The zero cusp value gives \(d\in W^t\).
The two generator equations then say
\[
               b-d\in W^s,\qquad c-d\in W^r .
 \tag{D.53.15}
\]
This is precisely (D.53.5). That map is
injective: an element in its kernel
would have \(d\in W^s\cap W^r=W^G=0\).
Taking dimensions proves (D.53.6).

\emph{The actual arithmetic centralizer data.}
The program
\texttt{d\_arithmetic\_correction\_graph.py}
enumerates the actual \(Q\) matrices modulo
twelve and checks the marked order \(576\).
It reconstructs all \(97\) free graph
cycles and both complete \(\mathbb F_3\)
action matrices. Exact elimination gives
\[
 \operatorname{rank}(s-1)=52,\qquad
 \operatorname{rank}(r-1)=63,\qquad
 \operatorname{rank}(r-1)^2=29 .
 \tag{D.53.16}
\]
These are the fixed dimension \(45\)
and the Jordan decomposition in (D.53.9).

The centralizer maps are checked in the
actual elementary-three extension, not
by enumerating a smaller substitute.
For each of the \(16\) or \(18\) matrices
centralizing \(s\) or \(r\) in \(Q\),
choose its recorded \(B\)-word.
Its commutator with the torsion generator
is an actual vector in \(N^{\rm ab}/3\).
A kernel correction centralizes the
generator exactly when
\[
              (1-g)v=-[\widetilde c,g]
                   \quad\hbox{in }N^{\rm ab}/3 .
 \tag{D.53.17}
\]
The certificate retains every such
equation and its solution or failure.
All \(16\) images lift for \(s\).
For \(r\), the liftable images are
exactly \(1,r,r^2\).
The two centralizer kernels have orders
\(3^{45}\) and \(3^{34}\), respectively.
This proves the exact centralizer
orders in (D.53.9).

For later use, every quotient of the
\(\mathbb F_3[\langle r\rangle]\)-module
\(N^{\rm ab}/3\) is generated by at most
\(29+5=34\) elements over that local
cyclic-group algebra. Its invariant
dimension equals its coinvariant
dimension, and is therefore at most
\(34\). Applying this to the actual
quotient by \(A\) gives
\[
             |C_{G/A}(\bar r)|\le18\cdot3^{34}.
 \tag{D.53.18}
\]
The factor \(18\) is the maximal
centralizer image in \(Q\); the kernel
is the just-bounded fixed subgroup.

\emph{Fixed dimensions and the positive lower bound.}
Write
\[
 q_s=|A^s|,\qquad a_t=|A\cap\langle t\rangle|,
 \qquad \bar c_r=|C_{G/A}(\bar r)|.
\]
The actual fixed dimensions are
\[
 \begin{aligned}
 \dim W^s&={w\over2}
               +4\cdot3^{45}(1-q_s^{-1}),\\
 \dim W^r&={w\over3}
               -{2\over9}(3^{35}-\bar c_r),\\
 \dim W^t&={n\over216}(1-a_t/q),
                    \qquad a_t\in\{1,3\}.
 \end{aligned}
 \tag{D.53.19}
\]
We justify each formula from its actual map.

Apply \(1-e_A\) to the graph sequence.
The augmentation disappears, leaving a
short exact edge--vertex sequence.
The moving edge module and the moving
\(r\)-vertex module are projective.
For the order-two subgroup, the periodic
resolution with maps \(1-s,1+s\) therefore
identifies \(H^1(\langle s\rangle,W)\)
with the degree-zero Tate group of the
moving \(s\)-vertex module.
A permutation module contributes one
such coordinate for each fixed vertex.
There are \(|C_G(s)|/2\) fixed vertices
before odd averaging, and
\(|C_{G/A}(\bar s)|/2\) afterwards.
The centralizer map onto the latter
group is onto: an involution discrepancy
in the odd elementary kernel is corrected
by solving \(1-s\) on its anti-invariant
part. Its kernel is \(A^s\).
Consequently
\[
 \dim H^1(\langle s\rangle,W)
    ={\,|C_G(s)|\over2}(1-q_s^{-1}).
\]
The left side is \(2\dim W^s-w\).
This gives the first formula.

For \(r\), odd-order invariants are exact.
There are \(n/3\) edge orbits, \(n/6\)
\(s\)-vertex orbits, and
\(n/9+2|C_G(r)|/9\) \(r\)-vertex orbits.
The last count uses
\(|C_G(r)|/3\) fixed \(r\)-vertices.
The classes of \(r,r^{-1}\) are distinct,
as is already checked in \(Q\).
The graph Euler formula thus gives
\[
       \dim V^r=1+n/18-2|C_G(r)|/9.
\]
The identical computation in \(G/A\)
computes \((V^A)^r\). Their difference
is the second formula in (D.53.19).

For the cusp, the certificate checks
that \(\langle t\rangle\) in \(Q\)
has order \(12\) and meets no conjugate
of either nontrivial elliptic subgroup.
It also checks that \(t^{12}\) is a
nonzero vector in \(N^{\rm ab}/3\);
the order in \(G\) is \(36\).
Thus \(a_t\) is one or three, and the
cusp order in \(G/A\) is \(36/a_t\).
Both cyclic subgroups act freely on
the vertices and edges of their graphs,
because any elliptic intersection
would already project to one in \(Q\).

The preimage in \(\widehat B\) of such
a cyclic subgroup of order \(h\) is
free of rank \(1+|G|/(6h)\).
The homology five-term sequence gives
the same dimension for \(V_{\langle t\rangle}\):
the free source has zero \(H_2\), and
the even cyclic quotient has one-dimensional
\(H_1\) and \(H_2\) over \(\mathbb F_2\).
For a cyclic operator, invariants and
coinvariants have equal dimensions.
Apply this first with \(h=36\) and then
with \(h=36/a_t\) in \(G/A\).
Subtracting gives the third formula.

Using (D.53.18) and \(a_t\ge1\), these formulas imply
\[
 \dim W^s\le w/2+4\cdot3^{45},\qquad
 \dim W^r\le w/3+4\cdot3^{34},\qquad
 \dim W^t\le w/36.
\]
Since \(q\ge3\), we also have \(w\ge n/9\).
Substitution in (D.53.6) gives
\[
 d\ge {5w\over36}-4\cdot3^{45}-4\cdot3^{34}
   \ge80\cdot3^{95}-4\cdot3^{45}-4\cdot3^{34}>0,
\]
as asserted.

\emph{Integral lifting and actual source automorphisms.}
The augmentation identity holds integrally:
\(\mathscr W_A=0\), because the odd
averaging idempotent is integral over
\(\mathbb Z_2\). Hence the map defining
\(\mathscr P\) is onto \(\mathscr W\).
It splits as a map of finite free
\(\mathbb Z_2\)-modules. Its kernel is
free of rank \(w\), and reduction gives
\[
                    \mathscr P/2\mathscr P
                          \xrightarrow{\sim}\mathcal P .
 \tag{D.53.20}
\]
This proves the required integrality
and reduction surjectivity.

D.50 identifies the actual kernel of
\(\mathcal Z_h\to\mathcal W_h\) with
the additive group \(\mathscr W\).
Left-multiply the two chosen conjugators
of \(\alpha_0\) by \(2b,2c\) in this
kernel. The old conjugates project to
\(s,r\) in \(G\), respectively; hence
their actions on this abelian kernel
are exactly the displayed \(s,r\)
module actions. Their exact cusp change is
\(2((1-s)b+s(1-r)c)=0\).
The torsion images remain conjugates.
On \(\mathcal W_h\) the action and
conjugators are unchanged. On the
complete homology component the
construction is an endomorphism,
since it sends \(K_3\) into its
abelian pro-two kernel and induces
identity on \(G\).

Write \(T_0\) for the invertible
integral homology matrix of \(\alpha_0\).
The new matrix is
\(T_0+2D_{b,c}\), where \(D_{b,c}\)
is the restriction of the displayed
integral derivation.
Its reduction modulo two is unchanged,
so it is invertible over \(\mathbb Z_2\).
It is therefore an automorphism of
the homology component. The unchanged
common-quotient action and literal
conjugator coordinates give an
automorphism on the actual fibre
product \(\mathcal Z_h\).
This proves the asserted integrated
family.

For \(\alpha_{2(b,c)}\alpha_0^{-1}\),
the matrix modulo four is
\[
                1+2\,\overline D_{b,c}\,
                                   \overline T_0^{-1}.
 \tag{D.53.21}
\]
All bars here denote reduction modulo two.
Multiplication of such matrices adds
their divided differences modulo two.
Their images are exactly the vector
space \(\mathscr D\,\overline T_0^{-1}\),
of dimension \(d\).
Thus the closed subgroup they generate
maps onto its additive elementary
abelian group, proving (D.53.11).
Every generator preserves the projection
to \(\mathcal J_\infty\), induces identity
on \(G\), and acts identically on \(U/2U\).
These properties persist under products,
inverses and closure. Consequently the
displayed modulo-four action is defined
on that whole closed subgroup, and
its kernel is precisely the elements
acting identically on \(U/4U\).
The map is continuous because it is
detected on the finite homology
quotient modulo four.
The entire D.49 action is retained;
its full loop abelianization also
retains the normalized commutator
parameter condition.

\emph{The separate finite same-action check.}
The program
\texttt{d\_cm\_same\_action\_obstruction.py}
rebuilds the order-\(1080\) marked graph
and computes the exact fixed spaces.
It obtains ranks \(56,82,120\) for
\(1-s,1-r,1-t\) on its moving module.
Thus (D.53.5) has kernel dimension
\(64+38=102\), with derivation image \(18\).
This checks explicitly that the
\(120\) pair parameters in D.52
were not \(120\) distinct derivations.

The full D.52 quadratic coefficients
are then evaluated on this kernel.
All kernel polar terms against all
\(120\) parameters vanish in the
recorded \(13\)-coordinate quotient.
The resulting same-action correction
image has rank \(4\). It contains the
actual obstruction of the earlier
pairing-preserving rank-\(16\) action.
An exact solution gives the new
conjugator parameter
\[
                  \mathtt{0x400000000000000000000002004840}.
 \tag{D.53.22}
\]
The program verifies equality of every
modulo-two homology column before and
after this change.

Finally
\texttt{d\_cm\_same\_action\_word.py}
uses that prescribed parameter in
the full D.52 word construction.
The entire \(4\,220\,100\)-letter
corrected cusp word is identity in
the unquotiented class-two target;
both torsion relations and rank-\(181\)
generation are checked.
The homology deviation remains rank
\(16\), and its matrix is exactly the
earlier pairing-preserving one.
The integral continuation through
exponents \(4,8,16\) retains this new
full class-two component.
These finite checks illustrate the
explicit local ambiguities. They do
not assign values to the arithmetic
invariant obstruction or lift the
automorphisms above to \(\widehat B\).
\end{proof}
